%! Setting Up Determinants --- the 2 by 2 Determinant, Area, and Orientation — Unit 5, Lesson 5.0 — Slides (Beamer)
\documentclass[aspectratio=169, 11pt]{beamer}
\usepackage{linalg-beamer}   % provides \burgundyheader{} and \sectionlabel[color]{}
\newcommand{\avec}{\mathbf{a}}
\newcommand{\bb}{\mathbf{b}}

\begin{document}

% TITLE SLIDE (CourseName is not defined in beamer, so the name is literal)
{
\setbeamercolor{background canvas}{bg=burgundy}
\begin{frame}[plain]
  \vspace{2.8cm}\hspace{0.55cm}%
  \begin{minipage}{0.9\textwidth}
    {\color{goldacc}\bfseries\small LESSON 5.0}\\[0.5cm]
    {\color{white}\bfseries\LARGE Setting Up Determinants --- the 2 by 2 Determinant, Area, and Orientation}\\[1.0cm]
    {\color{blushmid}\small Linear Algebra~$\cdot$~Unit 5: Determinants and Linear Transformations}
  \end{minipage}
\end{frame}
}

% HOOK
\begin{frame}[t, plain]
  \burgundyheader{Stretch a logo --- how much bigger, and did it flip?}
  \sectionlabel{Hook}
  \begin{columns}[T]
    \begin{column}{0.56\textwidth}
      Drop a logo in the unit square and stretch it with a matrix --- the square becomes a
      \textbf{parallelogram}.\\[8pt]
      \emph{One} number from the four entries answers both questions: how much \textbf{bigger}
      (the area) and whether it got \textbf{flipped} (the sign).
    \end{column}
    \begin{column}{0.42\textwidth}
      \centering
      \begin{tikzpicture}[scale=0.7]
        \draw[step=1, linegray!45, thin] (-0.3,-0.3) grid (4.4,2.4);
        \draw[->, thick, charcoal] (-0.4,0)--(4.6,0);
        \draw[->, thick, charcoal] (0,-0.4)--(0,2.6);
        \fill[burgundy!16] (0,0)--(3,0)--(4,2)--(1,2)--cycle;
        \draw[->, very thick, burgundy] (0,0)--(3,0);
        \draw[->, very thick, royalblue] (0,0)--(1,2);
        \draw[charcoal, thin] (3,0)--(4,2); \draw[charcoal, thin] (1,2)--(4,2);
      \end{tikzpicture}
    \end{column}
  \end{columns}
  \vspace{4pt}
  \begin{block}{The question of Unit 5}
    How can four numbers collapse into \textbf{one} that measures area and detects a flip?
  \end{block}
\end{frame}

% THE DETERMINANT
\begin{frame}[t, plain]
  \burgundyheader{The 2$\times$2 determinant}
  \sectionlabel[goldacc]{Definition}
  \begin{itemize}
    \setlength{\itemsep}{6pt}
    \item For $A=\begin{bsmallmatrix} a & b\\ c & d \end{bsmallmatrix}$, the \textbf{determinant} is
          \[ \det A = ad - bc \qquad(\text{down-diagonal} - \text{up-diagonal}). \]
    \item You already know it: it is the \emph{denominator} of the Unit~2 inverse
          $A^{-1}=\frac{1}{ad-bc}\begin{bsmallmatrix} d & -b\\ -c & a \end{bsmallmatrix}$.
    \item \textbf{Compute:} $\det\begin{bsmallmatrix} 2 & 1\\ 1 & 3 \end{bsmallmatrix}=6-1=5.$
  \end{itemize}
\end{frame}

% DETERMINANT = AREA
\begin{frame}[t, plain]
  \burgundyheader{The determinant is an area}
  \sectionlabel{Geometry}
  \begin{columns}[T]
    \begin{column}{0.54\textwidth}
      The two \textbf{columns} span a parallelogram; its \textbf{area} is $|\det A|$.\\[6pt]
      Columns $\begin{bsmallmatrix} 3\\0 \end{bsmallmatrix},\begin{bsmallmatrix} 1\\2 \end{bsmallmatrix}$:
      base $3$, height $2$, area $6$.\\[4pt]
      $\det\begin{bsmallmatrix} 3 & 1\\ 0 & 2 \end{bsmallmatrix}=6-0=6.$ \checkmark\\[6pt]
      Shearing slides the top --- base and height, hence area, stay the same.
    \end{column}
    \begin{column}{0.44\textwidth}
      \centering
      \begin{tikzpicture}[scale=0.8]
        \draw[step=1, linegray!45, thin] (-0.3,-0.3) grid (4.4,2.4);
        \draw[->, thick, charcoal] (-0.4,0)--(4.6,0);
        \draw[->, thick, charcoal] (0,-0.4)--(0,2.6);
        \fill[burgundy!16] (0,0)--(3,0)--(4,2)--(1,2)--cycle;
        \draw[->, very thick, burgundy] (0,0)--(3,0) node[below right]{\scriptsize $\begin{bsmallmatrix} 3\\0 \end{bsmallmatrix}$};
        \draw[->, very thick, royalblue] (0,0)--(1,2) node[above left]{\scriptsize $\begin{bsmallmatrix} 1\\2 \end{bsmallmatrix}$};
        \draw[charcoal, thin] (3,0)--(4,2); \draw[charcoal, thin] (1,2)--(4,2);
        \node[charcoal] at (2.2,1.0){\scriptsize area $=6$};
      \end{tikzpicture}
    \end{column}
  \end{columns}
\end{frame}

% ORIENTATION + ZERO
\begin{frame}[t, plain]
  \burgundyheader{The sign is orientation; zero means flat}
  \sectionlabel[goldacc]{Sign \& collapse}
  \begin{itemize}
    \setlength{\itemsep}{6pt}
    \item \textbf{Sign $=$ orientation:} $+$ keeps the turning direction, $-$ is a \emph{flip}.
          Swapping columns negates it: $\det\begin{bsmallmatrix} 1 & 3\\ 2 & 0 \end{bsmallmatrix}=-6.$
    \item \textbf{Zero $=$ flat:} parallel columns $\Rightarrow$ area $0$.
          $\det\begin{bsmallmatrix} 2 & 4\\ 1 & 2 \end{bsmallmatrix}=4-4=0.$
    \item \textbf{The whole story:} \; $\det A\ne0 \;\Leftrightarrow\;$ columns independent
          $\;\Leftrightarrow\; A$ is \textbf{invertible}.
  \end{itemize}
  \vspace{2pt}
  \begin{block}{The road ahead}
    \textbf{5.1} $3\times3$ determinants \& \emph{volume} \;$\bullet$\; \textbf{5.2} properties \&
    applications \;$\bullet$\; \textbf{5.3} transformations, where $|\det A|$ is the
    \emph{area-scaling factor}.
  \end{block}
\end{frame}

\end{document}
